% Contenido del póster: Braggion, Manconi, Pavanini y Zhu (2025), Journal of Financial Economics. DOI 10.1016/j.jfineco.2025.104113
% Elaborado con apoyo de IA a partir del resumen oficial del artículo (OpenAlex), 2026-11-12.
\def\DOI{10.1016/j.jfineco.2025.104113}
\def\TituloEs{El valor de la intermediación financiera: evidencia del crowdfunding de deuda en línea}
\def\TituloOriginal{The value of financial intermediation: Evidence from online debt crowdfunding}
\def\Autores{Braggion, Manconi, Pavanini y Zhu}
\def\Anio{2025}
\def\Revista{Journal of Financial Economics}

\def\Problema{La mayoría de los mercados en línea conectan directamente a las partes, pero los de crédito no: muchas plataformas de deuda han recuperado rasgos de los intermediarios financieros tradicionales. Cabe preguntarse por qué.}
\def\Objetivo{Usar el crédito en línea como laboratorio para medir el valor de la intermediación financiera.}
\def\Metodo{Los autores desarrollan un modelo estructural del crowdfunding de deuda en línea y lo estiman con una base de datos nueva. Con el modelo simulan escenarios alternativos, entre ellos uno en el que la plataforma funciona como un banco.}
\def\Resultados{\begin{itemize}
  \item Abandonar el esquema puro entre pares aumenta el excedente de los prestamistas, el beneficio de la plataforma y el crédito otorgado.
  \item A cambio, los inversionistas quedan expuestos al riesgo de liquidez.
  \item Si la plataforma asume ese riesgo, como un banco, el excedente y el crédito pueden ser aún mayores cuando la liquidez es escasa y los prestamistas tienen aversión al riesgo.
\end{itemize}}
\def\Conclusion{La intermediación tiene valor: las plataformas de deuda que adoptan funciones de intermediario benefician a prestamistas y prestatarios, aunque trasladan riesgos de liquidez a los inversionistas.}
\def\Relevancia{En el Perú, la plataforma líder opera con un balance pequeño frente a lo que intermedia: al cierre de 2025 tenía activos por S/~\cifra{activoLider2025} mil, mientras el mercado que lidera financió S/~\cifra{monto2025} millones ese año. El estudio ayuda a entender qué funciones de intermediario puede asumir una plataforma y qué riesgos quedan en manos de los inversionistas.}
\def\PalabrasClave{intermediación financiera, crowdfunding de deuda, modelo estructural, riesgo de liquidez, plataformas}
% --- Ficha del estudio (herramientas/posteres/ficha_articulos.py) ---
\def\Ambito{Plataforma de crowdfunding de deuda en línea}
\def\Enfoque{Modelo estructural estimado con datos nuevos}
\def\FiguraPeru{4}
% --- fin de la ficha ---
